\section{A written bit on the actual nonuniform detector pointer}
\label{sec:chvidr-detector-join}

The completed Chapter IV detector supplies an actual continuum measurement
and a finite uncentered second moment. Its positive intrinsic cap mean
remains a separate source-to-field transport operand. The construction
below uses the Chapter VI source decoder to write a calibrated displacement
onto that actual detector pointer. The unknown continuum signal is retained
in the background noise, rather than assigned a positive mean. Thus this
is an explicit added bit-to-meter receiving law, not a certificate that the
unmodified radiation alone distinguishes the two source ages.

The original sixteen source states, six signed lanes, both field
polarizations and counterterms, hard dressing, soft input and all angular
complements remain. One nominated lane and one of the actual local pointer
coordinates suffice for the binary receiving protocol. If both compatible
physical pointers are operated, the other pointer and its outcomes are
retained; no independence between their records is inferred.

\subsection{Re-key the field energy after the source decoder}
\label{sec:chvidr-postdecoder-energy}

Let the complete source-decoder output retain its two-state bit $B$ and
the source--field system. It can be correlated or entangled across these
recipients. In the selected reference lane use
\begin{equation}
\begin{gathered}
L=H_0-\epsilon eI+\epsilon I\geq0,
 \\
 H_f\leq2L+2cI,\\
 c=2g^2\mu.
\end{gathered}
\label{eq:chvidr-source-positive-energy}
\end{equation}
These are the original completed-square form inequalities; the scalar
subtraction is lane-specific. Suppose the reference decoder input has
$\langle L\rangle\leq C_{\rm dec}$ and its finite rectangular controls
have total absolute switching-work bound $W_{\rm dec}$. Their bounded
Hamiltonian jumps and conservation between jumps imply
\begin{equation}
\begin{gathered}
\langle L\rangle_{\rm out}\leq C_{\rm dec}+W_{\rm dec},
 \\
 E_{\rm rec}:=\langle H_f\rangle_{\rm out}
       \leq2(C_{\rm dec}+W_{\rm dec})+2c.
\end{gathered}
\label{eq:chvidr-rekey-field-energy}
\end{equation}
The nonselective source-bit observation, or merely retaining the coherent
bit, does not increase this source--field expectation. Do not condition
this energy bound on a correct decoder outcome. All bit outcomes remain
in the joint state. The old Chapter IV accepted-state energy and its
numerical noise table are not substituted after the new driven decoder.

For an actual preparation supplied only within trace distance
$\delta_{\rm prep}$ of that reference, apply the same entire receiving
channel to the two inputs and add $\delta_{\rm prep}$ once at the final
probability bound. No energy or form-domain assumption on the actual
preparation follows from trace distance alone. Equation~\eqref{eq:chvidr-rekey-field-energy}
and the ensuing second-moment calculation concern the finite-energy
reference input and its complete, unselected decoder output.

\subsection{The actual continuum gate on this new input}
\label{sec:chvidr-actual-pointer-input}

After the source bit is isolated, the actual nonuniform source--field
propagator and detector gate act trivially on $B$. For a prescribed
source-excluding local gate let $X$ be its actual integrated field
observable in the unmeasured joint Heisenberg picture. The original
retained reflection and causal commutator statements give the scalar
self-Magnus coefficient $c_h$. Write
\begin{equation}
\begin{gathered}
C_f\leq1.00301,\\
 c_{\max}=\frac{C_f^2V_c}{2},\\
 |c_h|\leq c_{\max}.
\end{gathered}
\label{eq:chvidr-pointer-magnus-bound}
\end{equation}
Here $\mu_c,V_c,D_c$ are the full compact-cap energy-dual, vacuum and
half-energy squared norms, and $\mu_{\rm sh}=1/(4R_s)$ is the original
source-shell norm. The source radius $R_s$ is distinct from the bit
register $B$. Let $t_*$ bound the actual new elapsed source-to-gate epoch,
including the declared receiving chronology. The backward smooth-test
estimate gives
\begin{equation}
 \begin{aligned}
 \langle X^2\rangle&\leq M_{\rm rec},\\
 M_{\rm rec}={}&6C_f^2\mu_c E_{\rm rec}+3C_f^2V_c\\
   &+48g^2C_f^2 n_+^2\mu_{\rm sh}t_*^2D_c,
 \qquad n_+=\frac{2003}{2000}.
 \end{aligned}
 \label{eq:chvidr-actual-field-second-moment}
\end{equation}
The input may be mixed and correlated with $B$: the field estimate uses
its unselected initial field energy and the unchanged bounded source-current
norms. The bit is a spectator in this evolution. No final old-Fock energy,
angular-vacuum factorization, Gaussian field state or positive cap mean
is inserted. The signed $g$ remains in the source law; its square appears
only in this moment bound.

Let $Q,P$ be the selected detector pointer with $[Q,P]=i$, and let
$kP^2$, $k>0$, be its declared kinetic Hamiltonian. The chosen $k$ retains
the Chapter IV field-form completion margin, or that margin is supplied
by the additional quadratic stabilizer below. The gate has duration
$\ell$ and normalized positive profile $h$. A small positive pointer mass
coefficient alone does not stabilize the bilinear field gate.

For clarity, make the gate's form certificate explicit. In the static
normal frame or the common moving frame used to construct the actual
joint propagator, let $F_t\geq0$ be the reference field-energy form and
suppose
\begin{equation}
\begin{gathered}
H_{\rm sf}(t)\geq a_0F_t-b_0I,\\
 a_0>0,
 \\
 \|a_t(f)\psi\|^2\leq\mu_*\,F_t[\psi].
\end{gathered}
\label{eq:chvidr-gate-form-certificate}
\end{equation}
The latter energy-dual bound is in the same frame, with
$\Phi_t(f)=a_t(f)+a_t(f)^*$. The finite coefficients $a_0,b_0,\mu_*$
belong to that actual compact-profile common-form certificate. They are
not a claim that its conjugating implementer preserves the old
$Q(H_f)$. If a fixed old-energy frame with the stated lower bound is
available, one may take $F_t=H_f$ and $\mu_*=\mu_c$; that substitution
must be certified, rather than inferred from the initial-energy bound.
On the product form domain, energy-dual Cauchy--Schwarz gives
\begin{equation}
 |h\langle P\Phi_t(f)\rangle|
 \leq\frac{a_0}{2}\langle F_t\rangle
       +\frac{2\mu_*}{a_0}h^2\langle P^2\rangle.
 \label{eq:chvidr-gate-young-bound}
\end{equation}
Thus add a known smooth compact pointer control $s(t)P^2$, with
\begin{equation}
 s(t)\geq\frac{2\mu_*}{a_0}h(t)^2.
 \label{eq:chvidr-pointer-stabilizer}
\end{equation}
The total gate form retains the lower bound
$(a_0/2)F_t+kP^2-b_0I$; a slightly larger $s$ gives a strict relative
form margin for the common-form construction. This is a declared
receiving actuator, not a change of the original source counterterms.
The compact actual frame supplies a finite common-form certificate;
numerical stabilizer resources require its actual coefficients.

For the retained receiving geometry there is a concrete specialization.
The original receiver radius is below $65$, the prescribed quarter-age
gate ends before $T_0$, and $65+n_+t_*<2\cdot10^6$. The retained compact
auxiliary is identically physical on that entire backward cap worldtube.
This is the causal-agreement statement used for the incoming homogeneous
transport estimate, so its exterior does not alter the local field gate.
Use this same fresh postdecoder stage, with its gate still before $T_0$;
the later translation and hold do not need to fit inside that interval.
The Chapter IV source-energy certificate has
$Q_\tau\geq(7/8)X$, $X\geq H_f/2$, and physical lapse at least
$n_-=1997/2000$. After the retained scalar shift its coefficient is at
least $7n_-/16>1/4$, so one may use $a_0=1/4$ with a strict margin.
Its inverse-frame weighted operator bound is
\begin{equation}
\begin{gathered}
\|\omega^{-1/2}(p\mathbin{\pm}q-I)\omega^{1/2}\|
        \leq401A,\\
 A=2^{-83}.
\end{gathered}
\label{eq:chvidr-detector-frame-dual}
\end{equation}
For a real coordinate cap test the pulled smearing is $(p\mathbin{\pm}q)f$.
Consequently $\mu_*\leq(1+401A)^2\mu_c<C_f^2\mu_c$, and a sufficient
choice is
\begin{equation}
 s(t)=8C_f^2\mu_c h(t)^2.
 \label{eq:chvidr-concrete-gate-stabilizer}
\end{equation}
This specialization uses the weighted inverse-frame bound, not an
assertion of old-Fock form-domain preservation. Source-region agreement
alone would be insufficient: the full detector-worldtube agreement is
used here explicitly. A different geometry or longer field gate requires
its own version of that causal and common-form certificate. Here
$S_g=8C_f^2\mu_c\int h(t)^2\,dt$, and the peak and total variation are
the same coefficient times $\|h\|_\infty^2$ and
$\operatorname{TV}(h^2)$, respectively.

Put $S_g=\int_{\rm gate}s(t)\,dt$. Since $P^2$ commutes with the field
gate interaction at every time, this addition changes only a known free
pointer shear. Its exact pointer equation is
\begin{equation}
 Q_{\rm gate}=Q+2(k\ell+S_g+c_h)P+X.
 \label{eq:chvidr-exact-gate-pointer}
\end{equation}
The initially prepared pointer is independent of the entire $B$--source--field
state. This pointer preparation is the only product preparation used in
the following Gaussian convolution; the bit and continuum need not be
independent.

\subsection{An exact finite controlled translation}
\label{sec:chvidr-controlled-translation}

Put $Z_B=|0\rangle\langle0|-|1\rangle\langle1|$. After the field gate
has ended, assign a local pointer-coordinate control of duration $\tau>0$,
\begin{equation}
 H_w(t)=kP^2+\lambda(t)Z_BP+\frac{\lambda(t)^2}{4k}I
       =k\left(P+\frac{\lambda(t)Z_B}{2k}\right)^2\geq0.
 \label{eq:chvidr-positive-bit-translation}
\end{equation}
The real pulse has area $\int_0^\tau\lambda(t)\,dt=a>0$.
All displayed generators commute, so
\begin{equation}
 U_w=\exp\!\left[-ik\tau P^2-iaZ_BP
                   -i\int_0^\tau\frac{\lambda(t)^2}{4k}\,dt\right].
 \label{eq:chvidr-exact-write-unitary}
\end{equation}
This is a direct sum of translated massive-pointer Hamiltonians on the
two bit sectors. The unbounded momentum is handled by its spectral
calculus and positive quadratic completion, rather than called a bounded
bit control. On constant pulse segments the operators are self-adjoint
on the momentum-squared domain; their form domain requires only finite
$\langle P^2\rangle$. The exponential is defined on the whole Hilbert
space. This adds a specified bit-to-meter interaction, not a derived
microscopic SAM actuator or an autonomous controller.

Focus the initial Gaussian pointer for the known combined gate-and-write
duration: set
\begin{equation}
\begin{gathered}
Y=Q+2[k(\ell+\tau)+S_g]P,\\
 \langle Y\rangle=\langle P\rangle=0,
 \\ \operatorname{Var}Y=\sigma^2,\quad
 \operatorname{Var}P=\frac1{4\sigma^2},\\
 \operatorname{Cov}_{\rm sym}(Y,P)=0.
\end{gathered}
\label{eq:chvidr-full-duration-focus}
\end{equation}
After writing and a further free hold $s\geq0$, the exact output coordinate
in the initial picture is
\begin{equation}
 Q_{\rm out}(s)=Y+X+2(c_h+ks)P+aZ_B.
 \label{eq:chvidr-written-pointer-output}
\end{equation}
The write drift and the stabilizer's known shear have been included in
the focus. Simply reusing a focus for the field gate alone would leave
additional drift in the noise. The stabilizer leaves the actual $X$,
$c_h$ and source backreaction law unchanged. The bit commutes with the actual joint field gate and with $X$,
even when its state is correlated with the field.

\subsection{A finite retention window without an intrinsic mean floor}
\label{sec:chvidr-pointer-retention}

Declare one read at any chosen hold time $s\in[0,T_A]$. The background
coordinate $Q_b(s)=Y+X+2(c_h+ks)P$ has uncentered second moment
\begin{equation}
 \langle Q_b(s)^2\rangle
   \leq M_{\rm rec}+\sigma^2+
                     \frac{(c_{\max}+kT_A)^2}{\sigma^2}.
 \label{eq:chvidr-held-background-bound}
\end{equation}
The cross expectations between the centered initial Gaussian variables
and $X$ vanish by their initial product preparation. This statement does
not factor an expectation involving $Z_B$ and $X$.
Choose the finite width
\begin{equation}
 \sigma^2=c_{\max}+kT_A>0.
\end{equation}
Then the following uniform bound holds:
\begin{equation}
 \begin{aligned}
 B_A&=M_{\rm rec}+2(c_{\max}+kT_A)\\
    &=6C_f^2\mu_c E_{\rm rec}+4C_f^2V_c
       \\ &\quad+48g^2C_f^2n_+^2\mu_{\rm sh}t_*^2D_c+2kT_A,\\
 \sup_{0\leq s\leq T_A}\langle Q_b(s)^2\rangle&\leq B_A.
 \end{aligned}
 \label{eq:chvidr-rekeyed-memory-noise}
\end{equation}
The positive width follows from the nonzero cap test $V_c>0$. Thus the
construction gives a finite noisy position record for a declared finite
hold interval. It does not assume an indefinitely stable free pointer or
an unmodelled sequence of nondestructive reads.

Decode pointer positions $q>0$ as bit zero and $q\leq0$ as bit one.
Since $Z_B$ and $Q_b(s)$ commute, their spectral probabilities are an
ordinary joint probability distribution, with arbitrary correlations.
If the pointer sign disagrees with the stored bit, then
$|Q_b(s)|\geq a$. Spectral Markov therefore gives
\begin{equation}
 P(\widehat r\ne r)
       \leq\frac{\langle Q_b(s)^2\rangle}{a^2}
       \leq\frac{B_A}{a^2}.
 \label{eq:chvidr-bit-transfer-error}
\end{equation}
No conditional moment divided by the probability of a correct source bit
is used. For each true source-age label $j$, if the source decoder has
error at most $\epsilon_{\rm src}^{\rm ref}$ on the reference input,
the union of the source-error and transfer-error events gives
\begin{equation}
 \epsilon_{\rm final}^{\rm ref}
   \leq\min\left\{1,\epsilon_{\rm src}^{\rm ref}+\frac{B_A}{a^2}\right\}.
 \label{eq:chvidr-reference-total-error}
\end{equation}
This inequality needs neither independence of errors nor a known mean of
the actual field. Any desired positive transfer allowance $\eta$ has the
finite sufficient displacement $a\geq\sqrt{B_A/\eta}$.

For a subsequent smooth sign-controlled memory write, let its trusted
positive and negative position bands start at $w$ and $-w$, with
$0<w<a$. The probability of being in the wrong band or its unresolved
central region obeys the same joint-event argument:
\begin{equation}
 P(\hbox{position outside its bit's trusted band})
       \leq\frac{B_A}{(a-w)^2}.
 \label{eq:chvidr-trusted-band-error}
\end{equation}
This is a one-time position probability at the specified memory-entry
time. Motion during a later memory gate requires that gate's own bound;
it is not an assertion that the pointer is parked.

If the signed displacement magnitude differs from $a$ by at most
$\delta_a$ and the threshold has position error at most $\delta_q$,
replace the denominator in Eq.~\eqref{eq:chvidr-bit-transfer-error} by
$(a-\delta_a-\delta_q)^2$, with a strictly positive margin. Nearest-bin
rounding of width $w$ can be included by increasing $\delta_q$ by $w/2$.
This is a finite threshold margin, not a demand for an infinitely fine
position ruler. A finite aperture $|q|\leq L_q$ can retain overflow as
an explicit failure outcome. Its additional probability is at most
\begin{equation}
 P({\rm overflow})\leq
       \frac{(\sqrt{B_A}+a+\delta_a)^2}{L_q^2}.
 \label{eq:chvidr-pointer-aperture-tail}
\end{equation}
The estimate uses the second-moment triangle inequality, with no
factorization of bit and background. No overflow is silently discarded.

\subsection{The complete receiving POVM and instrument}
\label{sec:chvidr-complete-instrument}

The actual Gaussian pointer variance at hold time $s$ is
$\nu(s)=\sigma^2+(c_h+ks)^2/\sigma^2$. Let $\Gamma_\nu$ be its centered
Gaussian probability measure and let $E_X$ be the actual spectral measure
of $X$. With $I_0=(0,\infty)$, $I_1=(-\infty,0]$ and
$\Pi_r=|r\rangle\langle r|_B$, the joint bit/output effects in the
unmeasured source--field picture are
\begin{equation}
\begin{gathered}
F_{r,y}=\Pi_r\int
     \Gamma_{\nu(s)}\bigl(I_y-x-a(-1)^r\bigr)\,E_X(dx),
 \\
 F_y=\sum_rF_{r,y}.
\end{gathered}
\label{eq:chvidr-translated-continuum-povm}
\end{equation}
These positive effects sum to the identity. They are the actual Chapter IV
Gaussian-convolution law with the added controlled displacement. They
retain the state-dependent continuum spectral law; it has not been
replaced by a fitted cap mean. Unknown $c_h$ has its stated uniform bound,
while its actual value defines the exact instrument.

For completeness let $U_{\rm all}(s)$ denote the actual nonuniform joint
evolution, field gate, write unitary and free hold, including the unmeasured
source--field drift. For the supplied independent Gaussian pointer state
$\sigma_D$, define
\begin{equation}
 \begin{aligned}
 \mathcal I_y(\rho)={}&(I\otimes\mathcal P_y)
     U_{\rm all}(s)(\rho\otimes\sigma_D)U_{\rm all}(s)^*
     (I\otimes\mathcal P_y),\\
 \mathcal P_y={}&\boldsymbol{1}_{I_y}(Q).
 \end{aligned}
 \label{eq:chvidr-retained-pointer-instrument}
\end{equation}
Here $\rho$ includes the bit, source and full field, and the pointer is
kept in the quantum output. The maps are completely positive and their
sum is trace preserving. A classical outcome register may be added as
$\sum_y|y\rangle\langle y|\otimes\mathcal I_y(\rho)$.
The overflow refinement likewise includes every outcome. Tracing some
recipients for a reported marginal does not remove them from the physical
resource account.

Equation~\eqref{eq:chvidr-retained-pointer-instrument} is the ideal sharp
position-event instrument. A discontinuous position projection generally
does not preserve finite pointer kinetic energy, so it is not assigned a
finite registration-work implementation. A smooth finite-energy target
instrument can instead use real smooth multipliers
\begin{equation}
\begin{gathered}
m_0(q)=\cos\vartheta(q),\\
 m_1(q)=\sin\vartheta(q),
 \\ \vartheta(q)=
 \begin{cases}\pi/2,&q\leq-w,\\0,&q\geq w,\end{cases}
 \\ 0\leq\vartheta\leq\pi/2,
\end{gathered}
\label{eq:chvidr-smooth-registration}
\end{equation}
with bounded derivatives and a smooth connection inside $(-w,w)$.
Replace each $\mathcal P_y$ in the retained-pointer map by $m_y(Q)$.
The resulting maps again form a complete instrument. Their effects on
the bit--source--field input are the convolution of $m_y^2$ with the
same displaced actual continuum Gaussian law. Outside the central band
they return the correct sign with certainty, so the bit-transfer error
is bounded by Eq.~\eqref{eq:chvidr-trusted-band-error}. For any finite
pointer momentum second moment, their nonselective output obeys
\begin{equation}
\begin{aligned}
 \sum_y\|P m_y(Q)\psi\|^2
 &=\|P\psi\|^2+\langle\psi,\vartheta'(Q)^2\psi\rangle\\
 &\leq\|P\psi\|^2+\|\vartheta'\|_\infty^2\|\psi\|^2.
\end{aligned}
\label{eq:chvidr-smooth-registration-energy}
\end{equation}
This identity extends to mixed and entangled inputs. It describes the
real-multiplier target instrument. The positive memory pulse in the next
section has the same effects but includes its coordinate-dependent common
phase: with $\vartheta=\pi A_w/2$, its Kraus operators are
$e^{-i\vartheta}\cos\vartheta$ and
$i e^{-i\vartheta}\sin\vartheta$. For that instrument the corresponding
identity is
$\sum_y\|PK_y\psi\|^2
=\|(P-\vartheta')\psi\|^2+\langle\vartheta'^2\rangle$.
Its finite-write estimate below includes this phase through the full
multiplier derivative bounds. A two-state register
has the bounded frozen-coordinate target dilation generated by
$\vartheta(Q)$ times its Pauli rotation. Its finite-duration physical
implementation must retain the active $kP^2$ drift and charge the
corresponding memory-gate error; the smooth target is not a hidden
pointer-parking assumption.

If the original preparation is within $\delta_{\rm prep}$ of the
reference and the specified access, storage or registration channel has
a further worst-label joint-state error budget $\delta_{\rm access}$,
contractivity of this same complete channel gives
\begin{equation}
 \epsilon_{\rm final}\leq\min\left\{1,
   \epsilon_{\rm src}^{\rm ref}+\frac{B_A}{a^2}
                  +\delta_{\rm prep}+\delta_{\rm access}\right\}.
 \label{eq:chvidr-all-outcome-error-budget}
\end{equation}
Use the calibrated denominator and overflow term when those outputs are
part of the declared protocol. The preparation distance is paid once on
the entire channel; it is not also included in
$\epsilon_{\rm src}^{\rm ref}$. If a previously quoted source ready-error
budget already includes that distance, reassemble its stated additive
components rather than adding the same preparation allowance twice.
Ready failures remain explicit, and counting a per-label failure at
most $f$ as unsuccessful gives the all-attempt error
$f+(1-f)\epsilon_{\rm final}$.

\subsection{Finite energy, switching work and causal event map}
\label{sec:chvidr-write-resources}

The focused Gaussian has finite pointer momentum energy
$k\langle P^2\rangle=k/(4\sigma^2)$ and finite initial position
variance $\sigma^2+[k(\ell+\tau)+S_g]^2/\sigma^2$. Momentum is conserved
through the field gate, translation and hold. For a rectangular write
with $\lambda=a/\tau$, the two jumps of
Eq.~\eqref{eq:chvidr-positive-bit-translation} have total absolute mean
work bounded by
\begin{equation}
 W_{\rm write}^{\rm abs}
     \leq2|\lambda|\sqrt{\langle P^2\rangle}
                   +\frac{\lambda^2}{2k}
     =\frac{a}{\tau\sigma}+\frac{a^2}{2k\tau^2}.
 \label{eq:chvidr-finite-write-work}
\end{equation}
This includes the positive scalar completion on both jumps. The signed
write payment telescopes for this commuting control; its absolute borrow,
controller capital and pointer preparation are still charged. The actual
field-gate switching, deposited energy and backreaction retain the Chapter IV
ledger. The added stabilizer has finite peak expected capital
$\|s\|_\infty/(4\sigma^2)$ and additional absolute work at most
\begin{equation}
 W_{\rm stab}^{\rm abs}\leq
       \langle P^2\rangle\int|s'(t)|\,dt
       =\frac{\operatorname{TV}(s)}{4\sigma^2},
 \label{eq:chvidr-stabilizer-work}
\end{equation}
including endpoint jumps if used. Its known focus preparation also has
the finite position variance just displayed. No global renormalized
vacuum-ramp work or thermal reset cost is
inferred from Eq.~\eqref{eq:chvidr-finite-write-work}.

One explicit finite model choice is
\begin{equation}\begin{gathered}
a=2^{10},\\
 k=2^{-40},\\
 \tau=2^{20},\\
 \lambda=2^{-10},\\
 T_A=2^{20}.
 \label{eq:chvidr-finite-translation-example}
\end{gathered}\end{equation}
Its pointer mass in the convention $k=1/(2M_D)$ is $2^{39}$, and its
write completion height is $\lambda^2/(4k)=2^{18}$. The focused variance
is $\sigma^2=c_{\max}+2^{-20}$ and
$B_A=B_{\rm N2}(E_{\rm rec},t_*)+2^{-19}$. With trusted-band width
$w=a/2$, Eq.~\eqref{eq:chvidr-trusted-band-error} is $4B_A/a^2$.
These numbers exhibit finite amplitude, capital and latency in model
units; they are not measured SI calibration. The field-gate stabilizer
still pays its actual certified $S_g$, peak and variation separately.

The source decoder, actual continuum gate and added translation have
separate finite durations. The retention window starts after writing;
the receiving result exists only after the chosen hold and registered
read. A spatially local $Z_BP$ actuator requires the bit carrier to be
co-located with the selected meter. If the decoded bit is delivered from
elsewhere, its causal transport, delay and possible errors belong in the
declared access map and the new $t_*$ chronology. It is not supplied
instantaneously to a remote detector. The same fixed schedule and
classification rule apply to both unknown source-age hypotheses.
The source-clock phase certificate is needed before the bit is decoded.
The new $t_*$ concerns the finite actual continuum stage through its field
gate; the long write and hold occur after that coupling is off and do not
extend a source-clock or field-gate certificate to later times.

This supplies a complete continuum-backed receiving instrument and a
finite noisy pointer-retention model under its stated actuator and access
laws. The intrinsic nonuniform cap-mean transfer, microscopic actuator,
achieved calibration and a physical carrier's durable-memory/access
performance remain distinct operands; none is replaced by the added
known displacement.

